\title{Large Language Models Can Automatically Engineer Features for Few-Shot Tabular Learning}

\begin{abstract}
Large Language Models (LLMs) have demonstrated exceptional capabilities in addressing complex and novel reasoning tasks, making them highly promising for tabular learning—a cornerstone of numerous practical domains. This work introduces an innovative in-context learning paradigm, \textsf{FeatLLM}, that leverages LLMs as feature constructors to transform the original dataset into representations that are well-optimized for tabular prediction tasks. The synthesized features are utilized by a straightforward downstream machine learning model (e.g., linear regression) to estimate class probabilities, achieving strong few-shot performance. Critically, \textsf{FeatLLM} relies exclusively on this basic predictive model with the derived features during inference. Distinct from prior LLM-based strategies, \textsf{FeatLLM} obviates the necessity of querying the LLM for every data instance during prediction. Additionally, it requires only API-based access to LLMs and is robust to prompt length constraints. Comprehensive evaluation over various tabular datasets spanning multiple domains illustrates that \textsf{FeatLLM} discovers effective rule-based features, exceeding the performance of alternative methods like TabLLM and STUNT by a substantial margin (an average improvement of 10\%).
\end{abstract}

\section{Introduction}
Large Language Models (LLMs) have, in recent times, exhibited remarkable prowess at crafting relevant outputs for novel tasks even in the absence of tailored fine-tuning~\cite{creswell2022selection,imani2023mathprompter,taylor2022galactica}. Their extensive pre-training on large-scale text sources imbues them with rich background knowledge, which can take the place of human expertise in numerous fields~\cite{Genomes2021,singhal2023large,wilcox2003role}, facilitating automation across a wide spectrum of use cases, from dialogue systems~\cite{finch2023leveraging} to automated code repair~\cite{fan2023automated}. Notably, equipping LLMs with a relevant task description and a handful of illustrative examples allows these models to contextualize new problems, drawing focus to salient features and informative instances~\cite{brown2020language,zhao2023survey}. Such behavior underscores their utility for data analysis and feature engineering.

LLMs' capacity to reason and incorporate prior knowledge has increasingly been harnessed for tabular learning challenges. For example, leveraging facts such as a correlation between aging and heightened disease risk can enhance predictive accuracy in medical scenarios. Recent techniques have involved translating task and feature information into natural language, reformatting the tabular input, and presenting this serialized data to the LLM~\cite{dinh2022lift,wang2023anypredict}. Subsequent model utilization typically entails either in-context prompting~\cite{nam2023semi} or efficient adaptation approaches such as parameter-efficient fine-tuning~\cite{dinh2022lift,hegselmann2023tabllm}. These strategies exploit the LLMs’ integrated knowledge, routinely surpassing traditional tabular models, particularly in low-data (few-shot) contexts~\cite{hegselmann2023tabllm}.

Existing LLM-based methods for tabular learning encounter several notable drawbacks. Firstly, they require at least one LLM inference per data point for end-to-end prediction, which incurs significant computational costs. Furthermore, achieving high prediction accuracy typically necessitates fine-tuning of the LLM, a process that is not feasible for LLMs lacking accessible training infrastructures or tuning APIs. This becomes particularly problematic with cutting-edge LLMs that restrict access solely to inference APIs~\cite{anil2023palm,openai2023gpt4}. Additionally, the majority of current strategies are ill-suited to managing extensive prompts~\cite{liu2023lost}. As the dimensionality of the tabular data increases and its textual representation expands, these approaches frequently become impractical or suffer degraded performance. Collectively, these issues introduce substantial obstacles for deploying LLM-based tabular learning systems in practical scenarios.

The method we introduce departs from the conventional paradigm of using LLMs exclusively for end-to-end predictions, instead positioning them as instruments for feature engineering that more effectively harness their existing knowledge. We present a new in-context learning method, \textsf{FeatLLM}, designed to decipher the “criteria” guiding the predictive mechanisms of LLMs. For example, when tasked with disease prediction, an LLM is capable of directly formulating and producing diagnostic rules based on relevant feature conditions. These derived criteria—or decision rules—are then leveraged to construct new features, which can replace the original ones in subsequent modeling tasks. This reallocation of the LLM’s capabilities allows for feature generation, enabling a simplified downstream process; once informative features have been identified, a complex machine learning model is no longer necessary during inference, thereby reducing latency. Taking advantage of LLMs' in-context learning properties, features may be elicited simply by providing curated example demonstrations within the prompts, circumventing the need for training procedures. Moreover, incorporating ensemble-inspired strategies~\cite{breiman2001random}, such as feature bagging, can alleviate the complications introduced by overly large prompt sizes.

\textsf{FeatLLM} organizes the prompt used for feature creation into two core phases: first, interpreting the underlying problem and deducing how features correlate with targets; second, leveraging insights gained along with a limited set of examples to devise clear decision rules that distinguish between classes. This explicit step-wise method enables LLMs to disregard irrelevant features when formulating rules. The resulting decision rules are applied to the dataset, treating them as programmatic expressions, and the data is converted into binary features that denote whether each sample conforms to each specified rule. Subsequently, a simple model is trained on these binary features to estimate class probabilities. The system’s reliability is further enhanced by adopting an ensemble mechanism, repeatedly generating features and using feature bagging. By controlling the number of features and samples included during bagging, \textsf{FeatLLM} avoids issues associated with overly large prompts. Since it operates independently of model training and incurs minimal inference costs, the LLM can be effectively deployed across a range of practical tasks.

We assess \textsf{FeatLLM} on 13 diverse tabular datasets under low-data settings and demonstrate its effectiveness and consistency. Our findings reveal that LLMs effectively harness both external domain knowledge and information inherent in the data itself to produce high-value features. Across all tested configurations, our approach surpasses several established few-shot learning methods. The implementation is accessible through an anonymized GitHub repository at~\texttt{https://github.com/Sungwon-Han/FeatLLM}.

\section{Related Work}

\subsection{Few-Shot Learning with Tabular Data} Few-shot learning has made it possible to derive transferable representations from datasets with limited labeling requirements~\cite{chen2018closer}. Although the initial emphasis was on image-based tasks, the methodology now exhibits robust generalization across diverse modalities, encompassing text, audio, and, more recently, tabular data~\cite{majumder2022few,nam2023stunt,schick2021exploiting}. The challenge of learning from small labeled sample sets lies in the tendency to capture incidental correlations~\cite{bartlett2021deep}. To address this, contemporary approaches have integrated additional informational resources or domain-tailored priors to instill favorable inductive biases during model development. One strategy includes leveraging large-scale, easily accessible unlabeled data together with informed augmentation to enable semi-supervised learning~\cite{ucar2021subtab,yoon2020vime}. Another avenue centers on task-independent unsupervised pre-training to obtain advantageous model initialization~\cite{somepalli2021saint}. These unsupervised tasks may involve techniques such as data masking or corruption, followed by reconstruction loss objectives~\cite{arik2021tabnet,majmundar2022met}, or the implementation of data augmentation for contrastive learning~\cite{bahri2022scarf,chen2020simple}. However, the intrinsic heterogeneity of tabular datasets complicates the search for universally relevant augmentations, with such tactics generally underperforming in few-shot scenarios~\cite{nam2023stunt}.

More recent developments have introduced training models across a spectrum of tabular datasets unrelated to the end task, with subsequent transfer learning applied to target applications~\cite{zhang2023generative,levin2022transfer,zhu2023xtab}. A representative example is TransTab~\cite{wang2022transtab}, which employs transformers to learn semantic relationships between table columns by exploiting column names rather than just cell values. Insights gleaned from varied tables and column contexts have enabled zero-shot and few-shot task inference. Alternatively, TabPFN~\cite{hollmann2022tabpfn} pre-trains a Bayesian neural network using synthetically generated data from mixed distributions. At inference, both training and test data from the target task are provided simultaneously without further training. The accumulated knowledge from diverse datasets has allowed these methods to surpass conventional tree-based algorithms and demonstrates marked utility for few-shot learning.

\subsection{Language-Interfaced Tabular Learning} Incorporating prior knowledge through large language models (LLMs)~\cite{anil2023palm,openai2023gpt4} presents a compelling direction. LLMs, with training grounded in vast and varied textual datasets, are known to adapt to novel tasks, reinforcing their value across multiple disciplines~\cite{brown2020language}. Recent investigations have sought to tap into the extensive prior knowledge contained in LLMs for tabular tasks. Typically, these methods convert tabular data to text format, integrating them as input prompts to LLMs~\cite{dinh2022lift,hegselmann2023tabllm,wang2023anypredict}. Providing tabular content, feature explanations, and task specifics within prompts enables inference regarding inter-feature and inter-task relationships. This line of work has developed toward either explicit model training via parameter-efficient tuning strategies such as LoRA~\cite{hu2021lora} or IA3~\cite{liu2022few}, or the use of in-context learning, where a handful of example tasks are included as demonstrations in prompts without further training~\cite{dinh2022lift,hegselmann2023tabllm,nam2023semi,slack2023tablet}.

While prior techniques have demonstrated notable success in advancing few-shot tabular learning, their dependence on routing inference through an LLM in every instance introduces practical limitations, particularly within industrial settings. Rather than adhering to the standard end-to-end, opaque LLM pipeline for prediction per sample, this work proposes an alternative approach that leverages the LLM to deduce class-specific conditions or rules. \textsf{FeatLLM} translates these inferred rules into executable code for generating additional features, thus facilitating inference without mandatory LLM interaction. It utilizes in-context learning to distill rules by drawing upon domain knowledge and information available from few-shot instances.

\section{Methods}
The primary objective of \textsf{FeatLLM} is to uncover the underlying ``rules"—those conditions that define each class—from the provided few-shot examples, as shown in Figure~\ref{fig:main_model}. After extracting these rules through an LLM, they are converted into binary feature representations for each sample, indicating the satisfaction of each rule. These binary features are subsequently processed using a dot product with non-negative trainable parameters to yield estimated class probabilities. Model training facilitates the evaluation of each feature's contribution to class prediction (Section~\ref{Sec:training}). The procedure utilizes bagging for repeated iterations, culminating in ensemble combination. Details regarding the problem setup and each procedural component are articulated below.

\vspace*{-0.12in}
\paragraph{Problem formulation.} We examine a tabular dataset containing $N$ labeled instances $\mathcal{D} = \{(\mathbf{x}^i, \mathbf{y}^i)\}_{i=1}^{N}$. This dataset $\mathcal{D}$ encompasses $d$ features, where each data point $\mathbf{x}^i$ is represented as a vector with $d$ dimensions. Associated with the features are descriptive names in natural language, defined by the set $F = \{f_j\}_{j=1}^{d}$, and accompanied by individual definitions. Under the classification paradigm, each label $\mathbf{y}^i$ is assigned from a set of designated categories. In $k$-shot learning scenarios, the training process utilizes only $k$ ($<N$) randomly chosen labeled instances from the dataset.

\subsection{Prompt Design for Extracting Rules}
\label{Sec:prompt_design}
To foster more precise rule extraction by the LLM, the prompting strategy emulates an expert's systematic approach to tabular data analysis. The designed prompt consists of three central segments (illustrated in Fig.~\ref{fig:prompt_for_rule} and detailed further in Appendix~\ref{sec:example_rule}):

\vspace*{-0.12in}
\paragraph{Basic information description.} This section outlines the foundational details needed to address the task (highlighted in orange within Fig.~\ref{fig:prompt_for_rule}). It presents a narrative describing the target problem, including information regarding possible labels, explanatory details for the features, and sample demonstrations. The narrative of the task is written as a question (for instance, ``Does this patient's myocardial infarction complications data indicate chronic heart failure? Yes or no?"), with additional task descriptions provided in Appendix~\ref{sec:task_desc}. Feature descriptions specify the type of value—possibly supplemented with definitions—and, for categorical variables, enumerate example classes where applicable. In the example demonstrations, a small subset of training examples is converted into textual format, paired with corresponding true labels. This textual serialization adheres to conventions established in prior work~\cite{dinh2022lift,hegselmann2023tabllm}:

\begin{align}
&\text{Serialize}(\mathbf{x}^i,\mathbf{y}^i, F) = \nonumber \\ 
&``\ f_1\ \text{is}\ \mathbf{x}^i_1.\ ... \ f_d\  \text{is}\  \mathbf{x}^i_d.\ \ \text{Answer:}\  \mathbf{y}^i\ ”
\end{align}

\vspace*{-0.12in}
\paragraph{Reasoning instruction.} To strengthen the reasoning abilities of the LLM, we incorporate reasoning instructions (highlighted in blue in Fig.~\ref{fig:prompt_for_rule}) into the prompting process. The reasoning instruction commences with an opening sentence inspired by the chain-of-thought methodology~\cite{wei2022chain}. Subsequently, we organize the LLM's reasoning workflow into two distinct stages. First, the LLM is prompted to hypothesize causal patterns or propensities between features and the task description, doing so independently of example demonstrations. This initial phase relies on general knowledge and common sense, enabling the LLM to utilize its prior understanding of the task. Next, in the second phase, the LLM integrates the provided example demonstrations with insights gained in the initial phase to generate a set of rules for each class. Employing this sequential reasoning approach mitigates the model's risk of recognizing spurious correlations from extraneous columns, thereby promoting attention toward relevant features. To balance the breadth and specificity of extracted rules—preventing both underfitting and overly verbose prompts—a fixed rule count (specifically, 10)\footnote{Analysis on the number of rules can be found in Section~\ref{sec:analysis}.} is imposed for each class. Further, we instruct the LLM to reference feature types as described in the basic information section during rule creation, reducing the likelihood of rule-feature type inconsistencies.

\vspace*{-0.12in}
\paragraph{Response instruction.} To ensure the inferred rules are straightforward to interpret and deploy, we include guidance on response formatting within the prompt (shown as yellow text in Fig.~\ref{fig:prompt_for_rule}). This prompt delineates both the appropriate answer format for each class and the expected structure of individual rules, helping the LLM choose rule formats compatible with the underlying feature type. With these directives, the LLM is empowered to logically combine rules—using operators such as AND or OR—when additional instructions are absent. An illustrative outcome is given in Appendix~\ref{sec:outcome-rule}.

\subsection{Inferring Class Likelihood via Rules}
\label{Sec:training}

\paragraph{Feature Generation via Parsing Rules.} We employ the previously generated rules to construct binary features for each class, specifying whether a given sample conforms to the rules corresponding to that class (resulting in either '0' or '1'). These binary indicators serve as inputs for estimating the probability that a sample pertains to a particular class. Since the LLM produces rules formulated in natural language, the automatic creation of these features necessitates an effective parsing method. Although we enforce structured formatting during rule generation to aid parsing, LLM outputs may still exhibit unwanted syntactic irregularities~\cite{kovavc2023large}. For example, the LLM may replace symbols such as ``$>$” or ``$<$” with expressions like ``greater than" or ``less than," complicating the parsing process.

To mitigate difficulties caused by such syntactic noise, rather than engineering elaborate parsing algorithms, we take advantage of the LLM's strengths. LLMs are adept at grasping semantic content regardless of superficial syntactic modifications and are proficient at composing straightforward code upon request (owing to their training on code repositories). To tap into these capabilities, we supply the LLM with prompts containing the function signature, explicit input/output specifications, and the extracted rules. These are then fed into the LLM for processing. Examples of this prompting approach and the resulting outputs are illustrated in Figures~\ref{fig:prompt_for_parsing} and~\ref{sec:example_parsing}-\ref{sec:example_parse_outcome} in the Appendix. The code generated by the LLM is subsequently executed via Python’s \texttt{exec()} function to facilitate feature extraction.

\vspace*{-0.12in}
\paragraph{Assessing Class Probabilities.} With binary rule-derived features for each class at hand, a straightforward technique for evaluating the likelihood that a sample belongs to a certain class is tallying the number of rules satisfied by the sample in each class (i.e., calculating the sum of that class’s binary features). Yet, since various rules may differ in significance, it is important to learn their respective contributions from the training data. To do so, we deploy a standard machine learning (ML) model tailored to the binary features for each class to quantify rule importance. For instance, consider using a linear model without bias. Define $\mathbf{z}_k^i \in \{0, 1\}^R$ as the vector of binary features produced for sample $\mathbf{x}^i$ for class $k$, where $R$ denotes the rule count for each class and $c$ is the total number of classes. The probability vector $\mathbf{p}^i$ for class predictions is determined as follows:

\begin{align}
\text{logit}_k^i &= \max(\mathbf{w}_k, 0) \cdot \mathbf{z}_k^i, \nonumber \\
\mathbf{p}^i &= \text{Softmax}({[\text{logit}_{1}^i, ..., \text{logit}_{c}^i]}). \label{eq:infer_prob}
\end{align}

In this context, $\mathbf{w}_k$ denotes the adjustable parameters within the linear model, indicating the contribution of each feature. By confining these weights to a positive subspace, we guarantee that the LLM-generated features do not decrease the predicted probability for any class during training. After generating the logit vector, the softmax function is employed to map logits to class probabilities. The cross-entropy loss serves as the objective function to learn $\mathbf{w}_k$, using a limited number of labeled examples for training.

\vspace*{-0.12in}
\paragraph{Ensembling with bagging.} In the final stage, we repeat the preceding steps multiple times, yielding several inference models. These models are collectively used in an ensemble to derive the final prediction. With a set of weight vectors $\{\mathbf{w}[t]\}_{t=1}^T$ obtained from $T$ iterations, we estimate the class probabilities by averaging the outputs $\mathbf{p}^i$, each calculated from the corresponding weight $\mathbf{w}[t]$ as outlined in Eq.~\ref{eq:infer_prob}.

To diversify the rule sets used for ensembling, we set a high temperature during LLM inference and shuffle the order of few-shot examples, as the arrangement of demonstrations can affect the LLM’s output~\cite{lu2022fantastically}\footnote{Further analysis regarding the diversity of generated rules is available in Appendix~\ref{sec:diversity}}. Bagging is implemented by sampling subsets of features or training instances in each trial, which is especially advantageous when there is a substantial number of features or samples. This ensemble strategy confers two major benefits. First, even when some rules produced by the LLM are erroneous in certain trials, correct rules from other runs can mitigate these errors, thus improving the system’s overall robustness. This approach leverages the self-consistency property of LLMs~\cite{wang2022self}, as rules that frequently appear across independent trials tend to be reliable. Second, ensemble bagging addresses prompt size constraints inherent to LLMs. If the input tabular data exceeds the model’s prompt limit, bagging facilitates manageable input size, helping sustain performance. Although an individual rule set may fall short of encompassing all feature interactions, an ensemble of several weaker models from different trials effectively offsets gaps in feature or instance coverage present in single trials.

\section{Experiments}
We assess the performance of \textsf{FeatLLM} using a range of tabular datasets, and carry out several analyses to better understand the mechanisms underlying our framework. Further experimental details—including studies on alternative LLM architectures, qualitative assessments, and additional findings—are provided in the Appendix due to constraints on manuscript length.

\subsection{Performance Evaluation}
\paragraph{Datasets.}
Our experimental setup involves 13 datasets designed for binary and multi-class classification tasks: (1) Adult~\cite{asuncion2007uci}—tasked with determining if an individual’s annual income surpasses \$50,000; (2) Bank~\cite{moro2014data}—identifies whether a client will opt for a term deposit; (3) Blood~\cite{yeh2009knowledge}—forecasts donor re-attendance for future donations; (4) Car~\cite{kadra2021well}—classifies automotive quality; (5) Communities~\cite{misc_communities_and_crime_183}—estimates regional crime rates; (6) Credit-g~\cite{kadra2021well}—evaluates creditworthiness as either favorable or unfavorable; (7) Diabetes\footnote{\texttt{kaggle.com/datasets/uciml/pima-indians-diabetes-database}}—predicts the likelihood of an individual having diabetes; (8) Heart\footnote{\texttt{kaggle.com/datasets/fedesoriano/heart-failure-prediction}}—detects coronary artery disease; and (9) Myocardial~\cite{misc_myocardial_infarction_complications_579}—predicts chronic heart failure in individuals.

Moreover, our analysis encompasses recent and synthetic datasets that are rarely discussed and are outside the typical scope of LLMs’ pre-training: (10) Cultivars—used to predict whether soybean varieties will yield high grain production\footnote{\texttt{archive.ics.uci.edu/dataset/913}}; (11) NHANES—classifies individuals’ age group from available health records\footnote{\texttt{archive.ics.uci.edu/dataset/887}}; (12) Sequence-type—categorizes sequences as one among arithmetic, geometric, Fibonacci, or Collatz; (13) Solution-mix—determines if a mixture from four solutions exceeds a concentration threshold of 0.5. The Cultivars and NHANES datasets, newly uploaded to the UCI Machine Learning Repository after September 2023, guarantee exclusion from the training data of the LLM API (gpt-3.5-turbo-0613) employed in this work. Sequence-type and Solution-mix are synthetically generated, ensuring the absence of memorization effects within the evaluated LLM.

These datasets exhibit diverse sizes and varying complexity, as reported in Table~\ref{Tab:basic-desc}.
Each dataset includes attributes that are labeled with descriptive names. High-dimensional datasets such as `Communities' and `Myocardial' contain over 100 attributes, introducing complexities related to the constrained context length of LLMs.

\vspace*{-0.12in}
\paragraph{Baselines.}
For comparative evaluation, we benchmark against nine different approaches. The initial three baselines are established tabular learning methods: (1) Logistic Regression (LogReg), (2) XGBoost~\cite{chen2016xgboost}, and (3) RandomForest~\cite{ho1995random}. Two subsequent baselines are geared toward leveraging unlabeled datasets: (4) SCARF~\cite{bahri2022scarf}, (5) STUNT~\cite{nam2023stunt}. However, acquiring vast collections of unlabeled samples may not always be practical. TabPFN~\cite{hollmann2022tabpfn} is a more recent baseline that relies on synthetic data generation with varied distributions to pretrain models. The last three baselines employ LLMs: (7) In-context learning (In-context)~\cite{wei2022emergent}, (8) TABLET~\cite{slack2023tablet}, and (9) TabLLM~\cite{hegselmann2023tabllm}. In-context learning involves embedding several example instances directly into the input prompt without fine-tuning. TABLET adds extra information to the prompt with supplementary rule-based or prototype hints generated by an external classifier, aiming to improve inference effectiveness. TabLLM adopts parameter-efficient adaptation approaches like IA3~\cite{liu2022few}, enabling the LLM to be trained using a few-shot paradigm.

\vspace*{-0.12in}
\paragraph{Implementation details.}
\textsf{FeatLLM} adopts GPT-3.5 as its foundational LLM, but its architecture remains independent of any particular LLM choice (see Appendix~\ref{sec:backbones} for empirical results using alternative backbones). The inference temperature is fixed at 0.5, while the top-p sampling parameter utilizes the default API value of 1. For feature extraction, the ensemble size and rule count are configured to 20 and 10, respectively. A comprehensive analysis of the hyper-parameter effects can be found in Figure~\ref{fig:hyperparameter2} and Appendix~\ref{sec:hyper-parameter}. Training of the linear model leverages the Adam optimizer with a learning rate of 0.01 over 200 epochs, where optimal epochs are selected by $k$-fold cross-validation. In baseline reproductions, GPT-3.5 is selected for in-context learning; meanwhile, T0~\cite{sanh2022multitask} is implemented for TabLLM, adhering to original specifications. Notably, TabLLM permits only LLMs with publicly accessible checkpoints, which precludes the use of GPT-3.5 in this scenario. For traditional methods like LogReg, XGBoost, and RandomForest, parameter tuning is accomplished via Grid-search with $k$-fold cross-validation. Here, $k$ is set to either 2 or 4 to ensure each class is represented in the training partition. For extensive implementation information, consult Appendix~\ref{sec:baselines}.

\vspace*{-0.12in}
\paragraph{Results.}
Tables~\ref{tab:main1} and \ref{tab:main2} provide comparative results of model performance across multiple datasets. Each experiment is executed three times with varying random splits of training and test data, and the reported AUC metrics include both their means and standard deviations. Across benchmarks, \textsf{FeatLLM} generally either leads or attains the second-best performance compared to baseline methods. Remarkably, our approach matches the results of standard tabular learning models utilizing 32 shots, even when only 4 shots are available (refer to Fig.~\ref{fig:compares}). Although STUNT and TabPFN excelled on select datasets, their cumulative advantage over traditional machine learning algorithms remained modest. In addition, large feature sets posed a challenge to LLM-driven methods such as in-context learning, TABLET, or TabLLM, which struggled with inference under these conditions. By incorporating bagging and ensemble methodologies, our framework handled high-dimensional datasets efficiently and achieved robust performance. It is worth mentioning that our bagging and ensemble framework can be integrated into other LLM-based techniques. Nevertheless, doing so would require doubling the inference steps per example, significantly heightening the computational load and rendering such adaptations less feasible in practice.

\subsection{Analysis \& Discussion}
\label{sec:analysis}
\paragraph{Ablation study.} We assess the role of principal framework components through targeted ablation studies: (1) \textsf{-Tuning} removes the weight adjustment phase from the linear module and simply aggregates binary feature values to determine the logit; (2) \textsf{-Ensemble} excludes the ensembling mechanism; (3) \textsf{-Description} omits the incorporation of feature descriptions; (4) \textsf{-Reasoning} removes Step-1 in the reasoning instruction stage during rule induction.

Table~\ref{tab:ablation} reports the average AUC impact of these ablations across datasets. All variations lead to some degree of diminished performance. Notably, eliminating weight tuning produces greater performance loss as the shot count increases, indicating that with more extensive data, precise rule weighting grows more influential for accuracy. Conversely, providing feature descriptions and structured reasoning yields substantial benefits under low-shot regimes, underscoring the importance of leveraging LLM's intrinsic prior knowledge for enhanced effectiveness in such scenarios. Finally, ensembling consistently boosts predictive outcomes in every experimental configuration.

\vspace*{-0.12in}
\paragraph{Training \& inference time comparison.} We evaluate the training and inference durations across multiple approaches. All experiments utilize the Adult dataset, with a training subset composed of 16 shots. Model execution is primarily conducted on a single A100 GPU, except for TabLLM, which leverages four A100 GPUs to facilitate model parallelism. Table~\ref{tab:cost} reports the runtime metrics. Remarkably, \textsf{FeatLLM} achieves inference times close to those of standard baselines, remaining efficient in comparison. Conversely, other large language model-based techniques, such as In-context learning, TABLET, and TabLLM, incur substantially greater inference times. For training, \textsf{FeatLLM} aligns with other LLM-based approaches, necessitating just 30 API calls, a figure that remains economically manageable and can be parallelized for increased efficiency.

\vspace*{-0.12in}
\paragraph{Quality of features generated by \textsf{FeatLLM}.}
To assess the informativeness of the rule-based features generated by \textsf{FeatLLM} for real-world applications, we carry out a straightforward experiment. For each dataset, the average AUROC between each newly constructed feature and the target label is computed, followed by determining the proportion of features exhibiting an AUROC above 0.5. The findings are presented in Table~\ref{tab:quality}. These results verify that most generated rules are pertinent to the prediction target and offer valuable information for the associated task (see the ``Before tuning'' column). Furthermore, during the linear model tuning stage, irrelevant rules—those with learned weights falling below a preset threshold—are automatically eliminated (see the ``After tuning'' column). The threshold is selected as 0.1, based on the overall weight histogram distribution.

\vspace*{-0.12in}
\paragraph{Impact of introducing spurious correlations.} To evaluate the capability of different approaches to mitigate the influence of extraneous spurious correlations in settings with limited data, we designed a targeted study. The Heart dataset was augmented by randomly incorporating columns from the Adult dataset, intentionally selecting features that do not pertain to the Heart dataset's classification objective. We then assessed how model performance shifted as we systematically increased the count of irrelevant columns. To ensure an equitable assessment, we fixed the number of labeled examples at 8 shots and did not leverage any unlabeled data, which omits approaches such as SCARF and STUNT from this evaluation.

The results, depicted in Figure~\ref{fig:spurious}, show the trajectory of performance as spurious correlations are introduced. In comparison to traditional baselines, \textsf{FeatLLM} exhibited the smallest decline in accuracy when confronted with spurious features. This robustness is attributable to the fact that our approach explicitly aligns rule extraction with prior knowledge concerning feature-task relationships, thereby discouraging the inclusion of rules based on irrelevant columns. Consequently, rules involving noisy features represented merely 1-2\% of the overall rule set. Nevertheless, it is also evident that \textsf{FeatLLM} may mistakenly identify and internalize spurious correlations, confusing the causal links between features and the target task, particularly when columns originating from semantically similar datasets are integrated at random (refer to Appendix~\ref{sec:spurious-appendix} for additional details).

\vspace*{-0.12in}
\paragraph{Hyper-parameter analysis}
\textsf{FeatLLM} is principally governed by two hyper-parameters: the ensemble count and the rule extraction count per class in each LLM invocation. To assess their influence on model efficacy, we ran empirical tests. Our findings indicate that increasing the ensemble number enhances performance, yet after about 20 ensembles, further improvements plateau (refer to Fig.~\ref{fig:appendix-hyperparameter1} in Appendix). Concerning the number of rules extracted per class, although variations did not yield significant performance differences, we selected 10 as the optimal setting, balancing prompt length considerations with predictive effectiveness (see Fig.~\ref{fig:hyperparameter2}). We hypothesize this is because adding more rules increases the input dimension, which provides richer information but also elevates the risk of overfitting, particularly in low-shot scenarios.

\section{Conclusion}
We introduce \textsf{FeatLLM}, a method that elevates the capabilities of LLM-driven few-shot tabular learning. Unlike approaches that simply leverage LLMs for per-instance prediction, \textsf{FeatLLM} specializes in generating predictive criteria akin to feature engineering. By employing a combination of rule creation and bagging-based ensemble techniques, \textsf{FeatLLM} minimizes inference overhead, requires only API usage (no fine-tuning), and circumvents feature dimension constraints. This framework achieves outstanding prediction results and represents a robust, efficient choice for deploying LLMs on practical tabular datasets.

\textsf{FeatLLM} is intentionally crafted for scenarios with limited data availability. Looking ahead, we intend to enhance its ability to construct new features for larger datasets, explore feature extraction strategies beyond rule-based forms, and advance interpretability—facilitating applicability across a broader spectrum of tasks.

\section{Broader Impact}
\textsf{FeatLLM} leverages the extensive prior knowledge embedded in LLMs for tabular data learning, aiming to surpass conventional supervised learning performances, particularly in data-constrained settings. By incorporating prior knowledge, our work effectively addresses overfitting challenges associated with restricted training examples, enhancing data efficiency. This approach proves advantageous in real-world domains such as finance and healthcare, where manual annotation is resource-intensive and LLMs trained on domain-related textual data can impart crucial supplementary insights. For future widespread adoption, it will be imperative to assess how societal biases or misinformation inherent in LLM pretraining may influence outcomes, especially given the possibility of such knowledge overshadowing observed labeled data.

\end{document}